\documentclass[11pt]{amsart}

% Encoding and fonts
\usepackage[T1]{fontenc}
\usepackage[utf8]{inputenc}
\usepackage{lmodern}

% Math packages
\usepackage{amsmath}
\usepackage{amssymb}
\usepackage{amsfonts}
\usepackage{amsthm}
\usepackage{mathtools}

% Layout and utilities
\usepackage{geometry}
\usepackage{enumitem}

\geometry{margin=1in}

\numberwithin{equation}{section}
\setcounter{tocdepth}{3}
\setcounter{secnumdepth}{3}

% ---------------------------------------------------------
% Theorem-like environments
% ---------------------------------------------------------

\theoremstyle{plain}
\newtheorem{theorem}{Theorem}[section]
\newtheorem{proposition}[theorem]{Proposition}
\newtheorem{lemma}[theorem]{Lemma}
\newtheorem{corollary}[theorem]{Corollary}
\newtheorem{conjecture}[theorem]{Conjecture}

\theoremstyle{definition}
\newtheorem{definition}[theorem]{Definition}
\newtheorem{experiment}[theorem]{Experiment}
\newtheorem{heuristic}[theorem]{Heuristic}
\newtheorem{question}[theorem]{Question}

\theoremstyle{remark}
\newtheorem{remark}[theorem]{Remark}
\newtheorem{warning}[theorem]{Warning}

\newcommand{\R}{\mathbb{R}}
\newcommand{\C}{\mathbb{C}}
\newcommand{\N}{\mathbb{N}}
\newcommand{\Qp}{\mathbb{Q}_{>0}}
\newcommand{\Hil}{\mathcal{H}}
\newcommand{\Sch}{\mathcal{S}}
\newcommand{\Four}{\mathcal{F}}
\newcommand{\Dom}{\mathcal{D}}
\newcommand{\Bop}{\mathcal{B}}
\newcommand{\id}{\mathrm{Id}}
\newcommand{\ad}{\operatorname{ad}}
\newcommand{\Reop}{\operatorname{Re}}
\newcommand{\Om}{\Omega}
\newcommand{\twconv}{\star_{\vartheta}}
\newcommand{\deltaone}{\delta_{1}}

% hyperref must be loaded last
\usepackage[hidelinks]{hyperref}

\title[Heisenberg-Twisted M\"obius Dynamics]{
Heisenberg-Twisted M\"obius Dynamics:\\
A Noncommutative Dirichlet Product, an Operator Zeta Inverse,\\
and Exact Position--Momentum Equations of Motion
}

\author{}
\date{}

\begin{document}

\maketitle

\tableofcontents

\bigskip

\begin{remark}[Mathematical status and novelty status]
\label{rem:status}
This is Version 2 of an exploratory research notebook.  The central change
from Version 1 is structural: the canonical commutation relation is no longer
placed next to an otherwise diagonal zeta multiplier.  Instead it deforms the
arithmetic product itself through the Weyl cocycle.  All operator domains,
convergence regions, and algebraic products used in the rigorous results are
specified below.

The Weyl representation, the canonical commutation relation, M\"obius
inversion, and the standard arithmetic functions used as ingredients are
standard.  The \emph{Heisenberg-twisted Dirichlet convolution}, the associated
M\"obius inverse coefficients $z_{\vartheta}(n)$, the prime-pair interference
formula in Theorem~\ref{thm:prime_pair_interference}, and the specific
M\"obius Hamiltonian dynamics in Section~\ref{sec:dynamics} are introduced
here as research-notebook constructions.  No claim of literature priority is
made in this version; a separate novelty audit would be required before any
priority statement.
\end{remark}

\section{Introduction: what Version 2 changes}
\label{sec:introduction}

The first version produced a mathematically valid Fourier-multiplier model,
but its Riemann and Heisenberg components were not dynamically coupled.  In
particular, a parameter $\sigma$ was external to the canonical pair, a
``Hamiltonian'' of the form $A^*A$ added no special dynamics, and plane waves
of exact momentum existed whether or not an arithmetic symbol vanished.  The
resulting non-localization statement was therefore a standard consequence of
$[Q,P]=i\hbar I$ appended to a separate zeta multiplier, rather than an
arithmetic consequence of the Heisenberg algebra.

Version 2 removes that architecture.  The starting point is the exact
Schr\"odinger pair
\begin{equation}
[Q,P]=i\hbar I,
\label{eq:intro_CCR}
\end{equation}
and its Weyl displacement operators.  An integer $n$ is assigned two
arithmetically defined phase-space increments:
\begin{equation}
x_n=\ell\log n,
\qquad
y_n=\kappa\Om(n),
\label{eq:intro_phase_space_coordinates}
\end{equation}
where $\Om(n)$ is the total number of prime factors of $n$, counted with
multiplicity.  The coordinate $x_n$ measures multiplicative size; on the
square-free support of $\mu$, the integer $\Om(n)$ is also the parity counter
in
\begin{equation}
\mu(n)=(-1)^{\Om(n)}.
\label{eq:intro_mobius_parity}
\end{equation}
Thus the two canonical directions encode two different arithmetic quantities
already present in M\"obius arithmetic.

The corresponding Weyl operators do not multiply according to the ordinary
commutative law of Dirichlet convolution.  The exact CCR produces a phase
cocycle.  This gives a new product $\twconv$ on arithmetic functions.  M\"obius
is then inverted with respect to this deformed product, producing coefficients
$z_{\vartheta}(n)$ which equal $1$ in the commuting limit but acquire explicit
cross-prime interference for nonzero Heisenberg area.

The central rigorous formula proved below is
\begin{equation}
z_{\vartheta}(pq)
=
2\cos\!\left(\frac{\vartheta}{2}\log\frac{q}{p}\right)-1
=
1-4\sin^2\!\left(\frac{\vartheta}{4}\log\frac{q}{p}\right)
\label{eq:intro_prime_pair_formula}
\end{equation}
for distinct primes $p,q$, whereas
\begin{equation}
z_{\vartheta}(p^r)=1
\label{eq:intro_prime_power_formula}
\end{equation}
for every prime power.  Equation~\eqref{eq:intro_prime_pair_formula} is not a
change of notation for the classical identity $\mu*1=\delta_1$: it is the
coefficient forced by inverting M\"obius in the noncommutative arithmetic
algebra generated by the Weyl form of Heisenberg's relation.

A genuine time evolution is then defined from the self-adjoint M\"obius
Hamiltonian $H_{\mu,\sigma,\vartheta}$.  Position and momentum are explicitly
$Q$ and $P$, and their exact Heisenberg equations are derived in
Theorem~\ref{thm:equations_of_motion}.  No variance or measurement-error
interpretation is required anywhere in that theorem.

\begin{warning}[What is not claimed]
\label{warn:no_RH_claim}
Version 2 does not claim a proof or a new equivalence of the Riemann
Hypothesis.  The ordinary zeta function is recovered at the commuting boundary
$\vartheta=0$.  For $\vartheta\neq0$ the inverse is a genuinely deformed zeta
object and its continuation and zero set are open problems.  The purpose of
this version is to produce a non-tautological coupling between Heisenberg
noncommutativity and M\"obius inversion before asking what, if anything, this
coupling implies about the classical zeros.
\end{warning}

\section{Exact canonical kinematics}
\label{sec:kinematics}

\subsection{Hilbert space, position, and momentum}
\label{subsec:QP_definitions}

\begin{definition}[Hilbert space]
\label{def:hilbert_space}
Fix $\hbar>0$ and set
\begin{equation}
\Hil=L^2(\R,dq)
\label{eq:hilbert_space}
\end{equation}
with its standard complex Hilbert-space inner product.
\end{definition}

\begin{definition}[Position]
\label{def:position}
The position operator $Q$ is
\begin{equation}
(Q\psi)(q)=q\psi(q)
\label{eq:Q_action}
\end{equation}
with maximal domain
\begin{equation}
\Dom(Q)=\{\psi\in L^2(\R):q\psi(q)\in L^2(\R)\}.
\label{eq:Q_domain}
\end{equation}
\end{definition}

\begin{definition}[Momentum]
\label{def:momentum}
The momentum operator $P$ is
\begin{equation}
(P\psi)(q)=-i\hbar\,\psi'(q)
\label{eq:P_action}
\end{equation}
with domain
\begin{equation}
\Dom(P)=H^1(\R).
\label{eq:P_domain}
\end{equation}
\end{definition}

\begin{proposition}[Self-adjointness and common core]
\label{prop:self_adjoint_QP}
The operators $Q$ and $P$ are self-adjoint.  The Schwartz space
$\Sch(\R)$ is a common invariant core for both.
\end{proposition}

\begin{proof}
$Q$ is the maximal multiplication operator by the real function $q$.  Under
the unitary Fourier transform with convention
\begin{equation}
(\Four\psi)(p)
=
\frac{1}{\sqrt{2\pi\hbar}}
\int_{\R}e^{-ipq/\hbar}\psi(q)\,dq,
\label{eq:fourier_transform}
\end{equation}
$P$ is unitarily equivalent to multiplication by the real variable $p$.
The core statement follows by standard cutoff and mollification arguments.
See, for example, \cite{ReedSimon1980}.
\end{proof}

\begin{proposition}[Exact canonical commutation relation]
\label{prop:CCR}
For every $\psi\in\Sch(\R)$,
\begin{equation}
[Q,P]\psi=i\hbar\psi.
\label{eq:CCR}
\end{equation}
\end{proposition}

\begin{proof}
Direct differentiation gives
\begin{align}
(QP\psi)(q)
&=-i\hbar q\psi'(q),
\label{eq:QP}\\
(PQ\psi)(q)
&=-i\hbar\psi(q)-i\hbar q\psi'(q).
\label{eq:PQ}
\end{align}
Subtracting \eqref{eq:PQ} from \eqref{eq:QP} proves \eqref{eq:CCR}.
\end{proof}

\subsection{Exact Heisenberg impossibility, without variances}
\label{subsec:heisenberg_impossibility}

\begin{theorem}[No simultaneous exact position and momentum]
\label{thm:no_joint_eigenstate}
There is no nonzero state with simultaneous exact position and exact momentum.
More precisely, both of the following are impossible for real $q_0,p_0$:
\begin{enumerate}[leftmargin=*]
\item a nonzero vector $\psi$ in a common invariant domain satisfying
\begin{equation}
Q\psi=q_0\psi,
\qquad
P\psi=p_0\psi;
\label{eq:joint_eigenstate_assumption}
\end{equation}
\item a nonzero tempered distribution $T\in\Sch'(\R)$ satisfying, in the
distributional sense,
\begin{equation}
qT=q_0T,
\qquad
-i\hbar T'=p_0T.
\label{eq:joint_distribution_assumption}
\end{equation}
\end{enumerate}
\end{theorem}

\begin{proof}
For vectors, \eqref{eq:joint_eigenstate_assumption} implies
\begin{equation}
[Q,P]\psi=0,
\label{eq:joint_eigenstate_commutator_zero}
\end{equation}
whereas \eqref{eq:CCR} gives
\begin{equation}
[Q,P]\psi=i\hbar\psi.
\label{eq:joint_eigenstate_commutator_CCR}
\end{equation}
Thus $\psi=0$.

For distributions define directly
\begin{equation}
QT=qT,
\qquad
PT=-i\hbar T'.
\label{eq:distribution_QP_definition}
\end{equation}
The distributional Leibniz rule gives
\begin{equation}
[Q,P]T=i\hbar T.
\label{eq:distribution_CCR}
\end{equation}
If \eqref{eq:joint_distribution_assumption} held, the left side would also
vanish, so $T=0$.
\end{proof}

\begin{remark}
\label{rem:no_measurement_uncertainty}
Theorem~\ref{thm:no_joint_eigenstate} is the exact algebraic content used in
this notebook when the word ``Heisenberg'' is invoked.  It is not a statement
about experimental error, measurement disturbance, or a product of standard
deviations.
\end{remark}

\section{Weyl displacements and the arithmetic phase space}
\label{sec:weyl_arithmetic}

\subsection{Weyl operators}
\label{subsec:weyl_operators}

\begin{definition}[Weyl displacement]
\label{def:weyl_operator}
For $(x,y)\in\R^2$, define
\begin{equation}
W(x,y)
=
\exp\!\left(\frac{i}{\hbar}(yQ-xP)\right).
\label{eq:weyl_definition}
\end{equation}
This is a unitary operator on all of $\Hil$.
\end{definition}

\begin{proposition}[Coordinate action]
\label{prop:weyl_coordinate_action}
For $\psi\in L^2(\R)$,
\begin{equation}
(W(x,y)\psi)(q)
=
\exp\!\left(\frac{i}{\hbar}y\left(q-\frac{x}{2}\right)\right)
\psi(q-x)
\label{eq:weyl_coordinate_action}
\end{equation}
almost everywhere.
\end{proposition}

\begin{proof}
Since $[iyQ/\hbar,-ixP/\hbar]$ is scalar, the Baker--Campbell--Hausdorff
formula gives
\begin{equation}
W(x,y)
=
e^{-ixy/(2\hbar)}e^{iyQ/\hbar}e^{-ixP/\hbar}.
\label{eq:weyl_BCH_factorization}
\end{equation}
The last factor translates by $x$ and the middle factor multiplies by
$e^{iyq/\hbar}$, yielding \eqref{eq:weyl_coordinate_action}.
\end{proof}

\begin{proposition}[Weyl form of the CCR]
\label{prop:weyl_relation}
Let $z=(x,y)$ and $z'=(x',y')$.  Define
\begin{equation}
\omega(z,z')=yx'-xy'.
\label{eq:symplectic_form}
\end{equation}
Then
\begin{equation}
W(z)W(z')
=
\exp\!\left(\frac{i}{2\hbar}\omega(z,z')\right)
W(z+z').
\label{eq:weyl_product}
\end{equation}
Consequently,
\begin{equation}
[W(z),W(z')]
=
2i\sin\!\left(\frac{\omega(z,z')}{2\hbar}\right)W(z+z').
\label{eq:weyl_commutator}
\end{equation}
\end{proposition}

\begin{proof}
For
\begin{equation}
A_z=\frac{i}{\hbar}(yQ-xP),
\label{eq:Az_definition}
\end{equation}
the CCR gives
\begin{equation}
[A_z,A_{z'}]
=
\frac{i}{\hbar}\omega(z,z')I.
\label{eq:Az_commutator}
\end{equation}
This commutator is central, so the Baker--Campbell--Hausdorff series stops
after the commutator term and gives \eqref{eq:weyl_product}.  Reversing $z$
and $z'$ changes the sign of $\omega$, giving \eqref{eq:weyl_commutator}.
The Weyl relations are standard; see \cite{Folland1989}.
\end{proof}

\begin{proposition}[Position and momentum displacement equations]
\label{prop:weyl_QP_commutators}
On $\Sch(\R)$,
\begin{align}
[Q,W(x,y)]
&=xW(x,y),
\label{eq:Q_W_commutator}\\
[P,W(x,y)]
&=yW(x,y).
\label{eq:P_W_commutator}
\end{align}
\end{proposition}

\begin{proof}
Both identities follow by differentiating \eqref{eq:weyl_coordinate_action},
or directly from \eqref{eq:CCR} because the relevant commutators with the
exponent are scalar.
\end{proof}

\subsection{Arithmetic embedding}
\label{subsec:arithmetic_embedding}

\begin{definition}[The additive prime-factor counter]
\label{def:Omega}
For
\begin{equation}
n=\prod_p p^{v_p(n)},
\label{eq:prime_factorization}
\end{equation}
define
\begin{equation}
\Om(n)=\sum_p v_p(n).
\label{eq:Omega_definition}
\end{equation}
Then
\begin{equation}
\Om(mn)=\Om(m)+\Om(n).
\label{eq:Omega_additivity}
\end{equation}
\end{definition}

\begin{definition}[Arithmetic phase-space map]
\label{def:arithmetic_phase_space_map}
Fix two positive scales $\ell>0$ and $\kappa>0$.  For $n\in\N$, put
\begin{align}
x_n
&=\ell\log n,
\label{eq:xn_definition}\\
y_n
&=\kappa\Om(n),
\label{eq:yn_definition}\\
z_n
&=(x_n,y_n).
\label{eq:zn_definition}
\end{align}
Define the dimensionless Heisenberg area parameter
\begin{equation}
\vartheta=\frac{\ell\kappa}{\hbar}.
\label{eq:theta_definition}
\end{equation}
\end{definition}

\begin{remark}[Why these two coordinates]
\label{rem:why_coordinates}
The first coordinate is forced by multiplicative size because
$\log(mn)=\log m+\log n$.  The second coordinate is chosen as the simplest
integer-valued completely additive function that directly controls the
M\"obius sign on square-free integers:
\begin{equation}
\mu(n)\neq0
\quad\Longrightarrow\quad
\mu(n)=(-1)^{\Om(n)}.
\label{eq:mobius_Omega_relation}
\end{equation}
The construction can be generalized by replacing $\Om$ with another
completely additive arithmetic function, but the present notebook fixes $\Om$ throughout.
\end{remark}

\begin{definition}[Arithmetic Weyl operators]
\label{def:arithmetic_weyl_operator}
For $n\in\N$, define
\begin{equation}
W_n=W(x_n,y_n).
\label{eq:Wn_definition}
\end{equation}
\end{definition}

\begin{theorem}[Exact Heisenberg cocycle on the integers]
\label{thm:integer_weyl_cocycle}
For $m,n\in\N$,
\begin{equation}
W_mW_n
=
\chi_{\vartheta}(m,n)W_{mn},
\label{eq:integer_weyl_product}
\end{equation}
where
\begin{equation}
\chi_{\vartheta}(m,n)
=
\exp\!\left[
\frac{i\vartheta}{2}
\bigl(\Om(m)\log n-\log m\,\Om(n)\bigr)
\right].
\label{eq:chi_definition}
\end{equation}
Moreover,
\begin{equation}
[W_m,W_n]
=
2i\sin\!\left[
\frac{\vartheta}{2}
\bigl(\Om(m)\log n-\log m\,\Om(n)\bigr)
\right]W_{mn}.
\label{eq:integer_weyl_commutator}
\end{equation}
\end{theorem}

\begin{proof}
Complete additivity gives
\begin{equation}
z_{mn}=z_m+z_n.
\label{eq:phase_space_additivity}
\end{equation}
Substitution of \eqref{eq:xn_definition}--\eqref{eq:theta_definition} into
\eqref{eq:weyl_product} gives \eqref{eq:integer_weyl_product};
\eqref{eq:integer_weyl_commutator} follows from
\eqref{eq:weyl_commutator}.
\end{proof}

\begin{corollary}[Prime commutator]
\label{cor:prime_commutator}
For distinct primes $p,q$,
\begin{equation}
[W_p,W_q]
=
2i\sin\!\left(\frac{\vartheta}{2}\log\frac{q}{p}\right)W_{pq}.
\label{eq:prime_commutator}
\end{equation}
Thus the prime operators commute only at the resonant values
\begin{equation}
\vartheta\log\frac{q}{p}\in2\pi\mathbb Z.
\label{eq:prime_commutation_resonance}
\end{equation}
\end{corollary}

\begin{remark}[Where Heisenberg enters nontrivially]
\label{rem:where_heisenberg_enters}
Equation~\eqref{eq:integer_weyl_product} is the first point where Version 2
ceases to be ordinary Dirichlet algebra in operator notation.  If the
canonical pair were commuting, the Weyl cocycle would be $1$ and the
arithmetic product below would reduce to ordinary Dirichlet convolution.
For nonzero $\vartheta$, the CCR creates order-dependent phases between
different arithmetic factors.
\end{remark}

\section{The Heisenberg-twisted Dirichlet algebra}
\label{sec:twisted_dirichlet_algebra}

\subsection{Definition and associativity}
\label{subsec:twisted_convolution_definition}

\begin{definition}[Heisenberg-twisted Dirichlet convolution]
\label{def:twisted_convolution}
For arithmetic functions $f,g:\N\to\C$, define
\begin{equation}
(f\twconv g)(k)
=
\sum_{d\mid k}
f(d)g(k/d)\chi_{\vartheta}(d,k/d).
\label{eq:twisted_convolution}
\end{equation}
The sum is finite for every $k$.
\end{definition}

\begin{lemma}[Cocycle identity]
\label{lem:cocycle_identity}
For all $a,b,c\in\N$,
\begin{equation}
\chi_{\vartheta}(a,b)\chi_{\vartheta}(ab,c)
=
\chi_{\vartheta}(b,c)\chi_{\vartheta}(a,bc).
\label{eq:cocycle_identity}
\end{equation}
\end{lemma}

\begin{proof}
The exponent in \eqref{eq:chi_definition} is the alternating bilinear form
obtained from the additive map $n\mapsto z_n$.  Expanding both sides using
$\log(ab)=\log a+\log b$ and
$\Om(ab)=\Om(a)+\Om(b)$ gives the same exponent.
\end{proof}

\begin{proposition}[Associative noncommutative arithmetic algebra]
\label{prop:twisted_associative}
The product $\twconv$ is associative and has identity $\deltaone$, where
\begin{equation}
\deltaone(1)=1,
\qquad
\deltaone(n)=0\quad(n>1).
\label{eq:delta_one_definition}
\end{equation}
For $\vartheta\neq0$ it is generally noncommutative.
\end{proposition}

\begin{proof}
Associativity follows from Lemma~\ref{lem:cocycle_identity} after grouping
the finite divisor sums by triples $abc=k$.  The identity property follows
from $\chi_{\vartheta}(1,n)=\chi_{\vartheta}(n,1)=1$.  Noncommutativity
follows from Corollary~\ref{cor:prime_commutator}; equivalently,
\begin{equation}
(\delta_p\twconv\delta_q)(pq)
=
\chi_{\vartheta}(p,q)
\label{eq:delta_pq_forward}
\end{equation}
whereas
\begin{equation}
(\delta_q\twconv\delta_p)(pq)
=
\chi_{\vartheta}(q,p)
=
\overline{\chi_{\vartheta}(p,q)}.
\label{eq:delta_pq_reverse}
\end{equation}
\end{proof}

\subsection{Weighted Banach algebra and operator representation}
\label{subsec:weighted_banach_algebra}

\begin{definition}[Weighted arithmetic Banach space]
\label{def:weighted_banach_algebra}
For $\sigma>0$, set
\begin{equation}
\mathcal A_{\sigma}
=
\left\{f:\N\to\C:
\|f\|_{\sigma}:=
\sum_{n\ge1}\frac{|f(n)|}{n^{\sigma}}<\infty
\right\}.
\label{eq:A_sigma_definition}
\end{equation}
\end{definition}

\begin{proposition}[Banach-algebra estimate]
\label{prop:banach_algebra_estimate}
For $f,g\in\mathcal A_{\sigma}$,
\begin{equation}
\|f\twconv g\|_{\sigma}
\le
\|f\|_{\sigma}\|g\|_{\sigma}.
\label{eq:banach_algebra_estimate}
\end{equation}
Hence $(\mathcal A_{\sigma},\twconv)$ is a unital Banach algebra.
\end{proposition}

\begin{proof}
Since $|\chi_{\vartheta}|=1$,
\begin{align}
\|f\twconv g\|_{\sigma}
&\le
\sum_{k\ge1}\frac{1}{k^{\sigma}}
\sum_{d\mid k}|f(d)||g(k/d)|
\label{eq:banach_estimate_step1}\\
&=
\sum_{d,e\ge1}
\frac{|f(d)|}{d^{\sigma}}
\frac{|g(e)|}{e^{\sigma}}
\label{eq:banach_estimate_step2}\\
&=
\|f\|_{\sigma}\|g\|_{\sigma}.
\label{eq:banach_estimate_step3}
\end{align}
Completeness is the completeness of weighted $\ell^1$.
\end{proof}

\begin{definition}[Weyl--Dirichlet representation]
\label{def:Weyl_Dirichlet_representation}
For $f\in\mathcal A_{\sigma}$, define
\begin{equation}
\Pi_{\sigma,\vartheta}(f)
=
\sum_{n\ge1}\frac{f(n)}{n^{\sigma}}W_n,
\label{eq:Pi_definition}
\end{equation}
where the series converges absolutely in operator norm.
\end{definition}

\begin{proposition}[Representation theorem]
\label{prop:representation_theorem}
For $f,g\in\mathcal A_{\sigma}$,
\begin{align}
\|\Pi_{\sigma,\vartheta}(f)\|
&\le\|f\|_{\sigma},
\label{eq:Pi_norm_bound}\\
\Pi_{\sigma,\vartheta}(f)
\Pi_{\sigma,\vartheta}(g)
&=
\Pi_{\sigma,\vartheta}(f\twconv g).
\label{eq:Pi_homomorphism}
\end{align}
\end{proposition}

\begin{proof}
The norm bound follows from $\|W_n\|=1$.  Absolute norm convergence permits
multiplication and regrouping of the double series.  The coefficient of
$W_k$ is exactly \eqref{eq:twisted_convolution} by
Theorem~\ref{thm:integer_weyl_cocycle}.
\end{proof}

\section{M\"obius inversion deformed by the CCR}
\label{sec:mobius_inversion}

\subsection{The M\"obius operator}
\label{subsec:mobius_operator}

\begin{definition}[M\"obius operator]
\label{def:mobius_operator}
For $\sigma>1$, define
\begin{equation}
\mathcal M_{\sigma,\vartheta}
=
\Pi_{\sigma,\vartheta}(\mu)
=
\sum_{n\ge1}
\frac{\mu(n)}{n^{\sigma}}W_n.
\label{eq:mobius_operator}
\end{equation}
The series converges absolutely in operator norm, so
$\mathcal M_{\sigma,\vartheta}\in\Bop(\Hil)$ and
\begin{equation}
\Dom(\mathcal M_{\sigma,\vartheta})=\Hil.
\label{eq:mobius_operator_domain}
\end{equation}
\end{definition}

\begin{remark}[Concrete coordinate action]
\label{rem:mobius_coordinate_action}
For $\psi\in L^2(\R)$,
\begin{equation}
(\mathcal M_{\sigma,\vartheta}\psi)(q)
=
\sum_{n\ge1}
\frac{\mu(n)}{n^{\sigma}}
\exp\!\left[
\frac{i}{\hbar}y_n\left(q-\frac{x_n}{2}\right)
\right]
\psi(q-x_n),
\label{eq:mobius_coordinate_action}
\end{equation}
with convergence in $L^2$.  Thus the operator is a M\"obius-weighted
superposition of simultaneous translations in position and momentum; it is
not a Fourier multiplier when $\kappa\neq0$.
\end{remark}

\subsection{Formal twisted inverse}
\label{subsec:formal_twisted_inverse}

\begin{proposition}[Existence and uniqueness of arithmetic inverses]
\label{prop:formal_inverse}
Let $f:\N\to\C$ satisfy $f(1)\neq0$.  Then there exists a unique arithmetic
function $g$ such that
\begin{equation}
f\twconv g=g\twconv f=\deltaone.
\label{eq:two_sided_inverse}
\end{equation}
Its coefficients are determined recursively from $g(1)=f(1)^{-1}$.
\end{proposition}

\begin{proof}
For $n>1$, the equation $(f\twconv g)(n)=0$ contains the term $f(1)g(n)$
and otherwise only values $g(d)$ with $d<n$.  This determines a unique left
inverse recursively.  The same recursion gives a unique right inverse.  In
an associative unital algebra, a left and right inverse coincide.
\end{proof}

\begin{definition}[Heisenberg-zeta coefficients]
\label{def:heisenberg_zeta_coefficients}
Let $z_{\vartheta}:\N\to\C$ be the unique two-sided inverse of $\mu$ for
$\twconv$:
\begin{equation}
\mu\twconv z_{\vartheta}
=
z_{\vartheta}\twconv\mu
=
\deltaone.
\label{eq:ztheta_inverse_definition}
\end{equation}
Thus
\begin{equation}
z_{\vartheta}(1)=1,
\label{eq:ztheta_one}
\end{equation}
and for $n>1$,
\begin{equation}
z_{\vartheta}(n)
=
-\sum_{\substack{d\mid n\\d>1}}
\mu(d)z_{\vartheta}(n/d)
\chi_{\vartheta}(d,n/d).
\label{eq:ztheta_recursion}
\end{equation}
\end{definition}

\begin{proposition}[Commuting boundary]
\label{prop:commuting_boundary}
At $\vartheta=0$,
\begin{equation}
z_0(n)=1
\label{eq:z0_constant_one}
\end{equation}
for every $n\ge1$.
\end{proposition}

\begin{proof}
When $\vartheta=0$, $\chi_0\equiv1$ and $\twconv$ is ordinary Dirichlet
convolution.  Equation~\eqref{eq:ztheta_inverse_definition} becomes
$\mu*z_0=\deltaone$.  Classical M\"obius inversion gives $z_0=\mathbf1$.
See \cite{Apostol1976}.
\end{proof}

\subsection{A bounded inverse in a rigorous half-plane}
\label{subsec:bounded_inverse}

\begin{theorem}[Uniform Neumann region]
\label{thm:neumann_region}
For every real $\vartheta$,
\begin{equation}
z_{\vartheta}\in\mathcal A_2.
\label{eq:ztheta_in_A2}
\end{equation}
Consequently, for every $\sigma\ge2$ the operator series
\begin{equation}
\mathcal Z_{\sigma,\vartheta}
=
\Pi_{\sigma,\vartheta}(z_{\vartheta})
=
\sum_{n\ge1}
\frac{z_{\vartheta}(n)}{n^{\sigma}}W_n
\label{eq:zeta_operator_deformed}
\end{equation}
converges absolutely in operator norm and satisfies
\begin{equation}
\mathcal M_{\sigma,\vartheta}
\mathcal Z_{\sigma,\vartheta}
=
\mathcal Z_{\sigma,\vartheta}
\mathcal M_{\sigma,\vartheta}
=I.
\label{eq:operator_inverse_identity}
\end{equation}
\end{theorem}

\begin{proof}
Write
\begin{equation}
\mu=\deltaone+h,
\qquad
h(1)=0,
\qquad
h(n)=\mu(n)\ (n>1).
\label{eq:mu_delta_h}
\end{equation}
Since $|\mu(n)|=\mu(n)^2$,
\begin{align}
\|h\|_2
&=
\sum_{n\ge2}\frac{\mu(n)^2}{n^2}
\label{eq:h_norm_step1}\\
&=
\frac{\zeta(2)}{\zeta(4)}-1
\label{eq:h_norm_step2}\\
&=
\frac{15}{\pi^2}-1
<1.
\label{eq:h_norm_step3}
\end{align}
The equality in \eqref{eq:h_norm_step2} is the standard Dirichlet series for
square-free integers.  Proposition~\ref{prop:banach_algebra_estimate} now
permits the Neumann series
\begin{equation}
(\deltaone+h)^{-1}
=
\sum_{r=0}^{\infty}(-h)^{\twconv r}
\label{eq:neumann_inverse}
\end{equation}
in $\mathcal A_2$.  By uniqueness of the arithmetic inverse, this inverse is
$z_{\vartheta}$.  Since $\mathcal A_2\subseteq\mathcal A_{\sigma}$ for
$\sigma\ge2$, Proposition~\ref{prop:representation_theorem} gives
\eqref{eq:operator_inverse_identity}.
\end{proof}

\begin{warning}[Exact border of the proved inverse]
\label{warn:inverse_border}
The M\"obius operator \eqref{eq:mobius_operator} exists as a bounded operator
for every $\sigma>1$.  The argument above proves invertibility uniformly in
$\vartheta$ only for $\sigma\ge2$.  No claim of invertibility for
$1<\sigma<2$ is made here.
\end{warning}

\section{The first Heisenberg arithmetic law}
\label{sec:prime_pair_interference}

\begin{theorem}[Prime rays are rigid; distinct primes interfere]
\label{thm:prime_pair_interference}
Let $p,q$ be distinct primes and let $r\ge1$.  Then
\begin{align}
z_{\vartheta}(p^r)
&=1,
\label{eq:ztheta_prime_power}\\
z_{\vartheta}(pq)
&=
2\cos\!\left(\frac{\vartheta}{2}\log\frac{q}{p}\right)-1,
\label{eq:ztheta_semiprime_cos}\\
z_{\vartheta}(pq)-1
&=
-4\sin^2\!\left(\frac{\vartheta}{4}\log\frac{q}{p}\right).
\label{eq:ztheta_semiprime_defect}
\end{align}
\end{theorem}

\begin{proof}
For $p^r$, only the divisors $1$ and $p$ contribute to the recursion because
$\mu(p^j)=0$ for $j\ge2$.  Moreover,
\begin{equation}
\Om(p)\log(p^{r-1})
-
\log p\,\Om(p^{r-1})
=0,
\label{eq:prime_power_symplectic_zero}
\end{equation}
so $\chi_{\vartheta}(p,p^{r-1})=1$.  Equation
\eqref{eq:ztheta_recursion} gives
\begin{equation}
z_{\vartheta}(p^r)=z_{\vartheta}(p^{r-1}),
\label{eq:prime_power_recursion}
\end{equation}
and \eqref{eq:ztheta_prime_power} follows from $z_{\vartheta}(1)=1$.

For $pq$, the four divisors $1,p,q,pq$ give
\begin{equation}
0
=
z_{\vartheta}(pq)
-
\chi_{\vartheta}(p,q)
-
\chi_{\vartheta}(q,p)
+1.
\label{eq:pq_inverse_equation}
\end{equation}
Since
\begin{equation}
\chi_{\vartheta}(p,q)
=
\exp\!\left(\frac{i\vartheta}{2}\log\frac{q}{p}\right),
\label{eq:chi_pq}
\end{equation}
and $\chi_{\vartheta}(q,p)$ is its complex conjugate,
\eqref{eq:ztheta_semiprime_cos} follows.  The identity
\eqref{eq:ztheta_semiprime_defect} is the elementary formula
$2(\cos u-1)=-4\sin^2(u/2)$.
\end{proof}

\begin{corollary}[Loss of ordinary multiplicativity]
\label{cor:loss_of_multiplicativity}
Generically,
\begin{equation}
z_{\vartheta}(pq)
\neq
z_{\vartheta}(p)z_{\vartheta}(q)=1.
\label{eq:loss_of_multiplicativity}
\end{equation}
Thus the Heisenberg deformation preserves every one-prime ray but destroys
ordinary multiplicativity through cross-prime interference.
\end{corollary}

\begin{corollary}[Small-area expansion]
\label{cor:small_area_expansion}
For fixed distinct primes $p,q$,
\begin{equation}
z_{\vartheta}(pq)
=
1-
\frac{\vartheta^2}{4}
\left(\log\frac{q}{p}\right)^2
+O(\vartheta^4)
\qquad(\vartheta\to0).
\label{eq:small_theta_pq}
\end{equation}
There is no linear correction.
\end{corollary}

\begin{proposition}[Evenness of the full inverse]
\label{prop:theta_evenness}
For every $n\ge1$,
\begin{equation}
z_{-\vartheta}(n)=z_{\vartheta}(n).
\label{eq:ztheta_evenness}
\end{equation}
\end{proposition}

\begin{proof}
Changing $\vartheta$ to $-\vartheta$ reverses the cocycle:
\begin{equation}
(f\star_{-\vartheta}g)(n)
=
(g\star_{\vartheta}f)(n).
\label{eq:opposite_algebra_identity}
\end{equation}
Since $z_{\vartheta}$ is a two-sided inverse of $\mu$ for
$\star_{\vartheta}$, it is a two-sided inverse of $\mu$ for the opposite
product $\star_{-\vartheta}$.  Uniqueness of the inverse gives
\eqref{eq:ztheta_evenness}.
\end{proof}

\begin{remark}[Why this is the non-tautological step]
\label{rem:non_tautological_step}
At $\vartheta=0$ the inverse coefficients are identically $1$.  At nonzero
$\vartheta$, equation~\eqref{eq:ztheta_semiprime_defect} creates a new
arithmetic coefficient depending on the logarithmic separation of two
primes.  This dependence is forced by the symplectic area in the Weyl
product and disappears if the CCR is removed.  It is therefore the first
result in the notebook where Heisenberg noncommutativity changes the
arithmetic inverse itself rather than merely constraining an unrelated
state space.
\end{remark}

\section{The Heisenberg-deformed zeta series}
\label{sec:deformed_zeta_series}

\begin{definition}[Heisenberg-deformed zeta series]
\label{def:deformed_zeta_series}
For $\Re s\ge2$, define
\begin{equation}
\zeta_{\mathrm H}(s;\vartheta)
=
\sum_{n\ge1}\frac{z_{\vartheta}(n)}{n^s}.
\label{eq:deformed_zeta_series}
\end{equation}
Theorem~\ref{thm:neumann_region} guarantees absolute convergence on the
closed half-plane $\Re s\ge2$.
\end{definition}

\begin{proposition}[Classical boundary]
\label{prop:zeta_classical_boundary}
For $\Re s>1$,
\begin{equation}
\zeta_{\mathrm H}(s;0)=\zeta(s).
\label{eq:zeta_classical_boundary}
\end{equation}
\end{proposition}

\begin{proof}
By Proposition~\ref{prop:commuting_boundary}, $z_0(n)=1$ for all $n$.
Substitution into \eqref{eq:deformed_zeta_series} gives the ordinary
Dirichlet series for $\zeta(s)$.
\end{proof}

\begin{proposition}[Even Heisenberg deformation]
\label{prop:zeta_theta_even}
For $\Re s\ge2$,
\begin{equation}
\zeta_{\mathrm H}(s;-\vartheta)
=
\zeta_{\mathrm H}(s;\vartheta).
\label{eq:zeta_theta_even}
\end{equation}
\end{proposition}

\begin{proof}
This follows termwise from Proposition~\ref{prop:theta_evenness} and absolute
convergence.
\end{proof}

\begin{warning}[No analytic continuation is asserted]
\label{warn:no_continuation}
The definition \eqref{eq:deformed_zeta_series} is presently rigorous only in
the half-plane furnished by Theorem~\ref{thm:neumann_region}.  No meromorphic
continuation, functional equation, Euler product, or critical-strip theory
for $\zeta_{\mathrm H}(s;\vartheta)$ is asserted here.
\end{warning}

\begin{proposition}[Conditional first-order rigidity of simple zeros]
\label{prop:conditional_zero_rigidity}
Assume that, for some simple zero $\rho$ of the ordinary zeta function,
there exists a function $F(s,\vartheta)$ analytic in a neighbourhood of
$(\rho,0)$ such that
\begin{align}
F(s,0)
&=\zeta(s),
\label{eq:F_classical_boundary}\\
F(s,-\vartheta)
&=F(s,\vartheta),
\label{eq:F_evenness}
\end{align}
and that $F$ agrees with $\zeta_{\mathrm H}$ wherever both are defined.
Then the local zero branch $\rho(\vartheta)$ through $\rho$ satisfies
\begin{equation}
\rho'(0)=0.
\label{eq:zero_first_derivative_zero}
\end{equation}
\end{proposition}

\begin{proof}
The implicit-function theorem applies because $\zeta'(\rho)\neq0$.
Evenness \eqref{eq:F_evenness} makes the unique local zero branch even, so
its first derivative at the origin vanishes.
\end{proof}

\begin{warning}
\label{warn:conditional_zero_result}
Proposition~\ref{prop:conditional_zero_rigidity} is conditional on an analytic
continuation in both variables and does not supply such a continuation.  It
is included because it identifies a precise way in which the exact
$\vartheta\mapsto-\vartheta$ symmetry would constrain the motion of classical
Riemann zeros under the Heisenberg deformation if that continuation exists.
\end{warning}

\section{A genuine M\"obius Hamiltonian}
\label{sec:dynamics}

\subsection{Definition and domain}
\label{subsec:hamiltonian_definition}

\begin{definition}[Cosine and sine Weyl quadratures]
\label{def:Weyl_quadratures}
For $n\ge1$, set
\begin{align}
C_n
&=\frac{W_n+W_n^*}{2},
\label{eq:Cn_definition}\\
J_n
&=\frac{W_n-W_n^*}{2i}.
\label{eq:Jn_definition}
\end{align}
Both are bounded self-adjoint operators with norm at most $1$.
\end{definition}

\begin{definition}[M\"obius Hamiltonian]
\label{def:mobius_hamiltonian}
Fix an energy scale $E_0>0$ and $\sigma>1$.  Define
\begin{equation}
H_{\mu,\sigma,\vartheta}
=
E_0\left(\Reop\mathcal M_{\sigma,\vartheta}-I\right)
=
E_0\sum_{n\ge2}
\frac{\mu(n)}{n^{\sigma}}C_n.
\label{eq:mobius_hamiltonian}
\end{equation}
The series converges absolutely in operator norm.
\end{definition}

\begin{proposition}[Self-adjoint bounded dynamics]
\label{prop:hamiltonian_selfadjoint}
For $\sigma>1$, $H_{\mu,\sigma,\vartheta}$ is bounded and self-adjoint on
all of $\Hil$:
\begin{equation}
\Dom(H_{\mu,\sigma,\vartheta})=\Hil.
\label{eq:hamiltonian_domain}
\end{equation}
It therefore generates the unitary group
\begin{equation}
U(t)=\exp\!\left(-\frac{i}{\hbar}tH_{\mu,\sigma,\vartheta}\right).
\label{eq:time_evolution_unitary}
\end{equation}
\end{proposition}

\begin{proof}
The norm series in \eqref{eq:mobius_hamiltonian} is dominated by
$E_0\sum_{n\ge2}|\mu(n)|n^{-\sigma}$, which converges for $\sigma>1$.
Every $C_n$ is self-adjoint, hence so is the norm limit.  The spectral theorem
then defines \eqref{eq:time_evolution_unitary}.
\end{proof}

\begin{proposition}[Explicit Schr\"odinger equation]
\label{prop:schrodinger_equation}
For $\psi_0\in\Hil$, the evolution $\psi(t)=U(t)\psi_0$ satisfies
\begin{equation}
i\hbar\,\partial_t\psi(t,q)
=
(H_{\mu,\sigma,\vartheta}\psi(t))(q)
\label{eq:schrodinger_equation}
\end{equation}
in the strong $L^2$ sense, where
\begin{align}
(H_{\mu,\sigma,\vartheta}\psi)(q)
=
\frac{E_0}{2}
\sum_{n\ge2}
\frac{\mu(n)}{n^{\sigma}}
\Bigg[
& e^{\frac{i}{\hbar}y_n(q-x_n/2)}\psi(q-x_n)
\notag\\
+{}&e^{-\frac{i}{\hbar}y_n(q+x_n/2)}\psi(q+x_n)
\Bigg].
\label{eq:hamiltonian_coordinate_action}
\end{align}
\end{proposition}

\begin{proof}
Equation~\eqref{eq:hamiltonian_coordinate_action} follows from
Proposition~\ref{prop:weyl_coordinate_action} applied to $W_n$ and $W_n^*$.
The series converges in operator norm, so the bounded Hamiltonian gives a
strong Schr\"odinger evolution for every initial vector.
\end{proof}

\subsection{Exact equations of motion for position and momentum}
\label{subsec:equations_of_motion}

\begin{lemma}[Bounded first commutators]
\label{lem:bounded_first_commutators}
For $\sigma>1$, the series
\begin{align}
\sum_{n\ge2}
\frac{|\mu(n)|x_n}{n^{\sigma}},
\label{eq:x_commutator_sum}\\
\sum_{n\ge2}
\frac{|\mu(n)|y_n}{n^{\sigma}}
\label{eq:y_commutator_sum}
\end{align}
converge.  Hence $[H_{\mu,\sigma,\vartheta},Q]$ and
$[H_{\mu,\sigma,\vartheta},P]$, initially defined on $\Sch(\R)$, extend to
bounded operators.
\end{lemma}

\begin{proof}
The first series is bounded by a constant multiple of
$\sum n^{-\sigma}\log n$.  Since
\begin{equation}
\Om(n)\le\frac{\log n}{\log2},
\label{eq:Omega_log_bound}
\end{equation}
the second is bounded by another constant multiple of the same convergent
series.
\end{proof}

\begin{theorem}[Exact Heisenberg equations of motion]
\label{thm:equations_of_motion}
Let
\begin{align}
Q(t)
&=U(t)^*QU(t),
\label{eq:Qt_definition}\\
P(t)
&=U(t)^*PU(t).
\label{eq:Pt_definition}
\end{align}
On the common invariant core, and equivalently in the associated quadratic
forms, the Heisenberg equations are
\begin{align}
\dot Q(t)
&=
\frac{E_0}{\hbar}
U(t)^*
\left[
\sum_{n\ge2}
\frac{\mu(n)x_n}{n^{\sigma}}J_n
\right]
U(t),
\label{eq:Q_equation_motion}\\
\dot P(t)
&=
\frac{E_0}{\hbar}
U(t)^*
\left[
\sum_{n\ge2}
\frac{\mu(n)y_n}{n^{\sigma}}J_n
\right]
U(t).
\label{eq:P_equation_motion}
\end{align}
At $t=0$ this becomes
\begin{align}
\dot Q(0)
&=
\frac{E_0\ell}{\hbar}
\sum_{n\ge2}
\frac{\mu(n)\log n}{n^{\sigma}}J_n,
\label{eq:Qdot_zero}\\
\dot P(0)
&=
\frac{E_0\kappa}{\hbar}
\sum_{n\ge2}
\frac{\mu(n)\Om(n)}{n^{\sigma}}J_n.
\label{eq:Pdot_zero}
\end{align}
\end{theorem}

\begin{proof}
The time-independent Heisenberg equation is
\begin{equation}
\dot A(t)
=
\frac{i}{\hbar}U(t)^*[H_{\mu,\sigma,\vartheta},A]U(t).
\label{eq:general_heisenberg_equation}
\end{equation}
From Proposition~\ref{prop:weyl_QP_commutators},
\begin{align}
[Q,C_n]
&=ix_nJ_n,
\label{eq:Q_Cn_commutator}\\
[P,C_n]
&=iy_nJ_n.
\label{eq:P_Cn_commutator}
\end{align}
Therefore
\begin{align}
[H_{\mu,\sigma,\vartheta},Q]
&=
-iE_0\sum_{n\ge2}
\frac{\mu(n)x_n}{n^{\sigma}}J_n,
\label{eq:H_Q_commutator}\\
[H_{\mu,\sigma,\vartheta},P]
&=
-iE_0\sum_{n\ge2}
\frac{\mu(n)y_n}{n^{\sigma}}J_n.
\label{eq:H_P_commutator}
\end{align}
Lemma~\ref{lem:bounded_first_commutators} justifies the sums as bounded
operators.  Substitution into \eqref{eq:general_heisenberg_equation} gives
\eqref{eq:Q_equation_motion} and \eqref{eq:P_equation_motion}.
\end{proof}

\begin{remark}[Physical and arithmetic meaning of the two equations]
\label{rem:meaning_equations_motion}
The position variable is $Q$, the logarithmic-scale coordinate.  The
momentum variable is $P=-i\hbar\partial_q$, the generator conjugate to that
coordinate.  Multiplication by an integer $n$ is represented by a phase-space
displacement of size $\ell\log n$ in $Q$ and $\kappa\Om(n)$ in $P$.
Equations~\eqref{eq:Qdot_zero} and \eqref{eq:Pdot_zero} therefore say that
the M\"obius Hamiltonian drives position by logarithmic size and momentum by
prime-factor count, with the same quantum current operators $J_n$.
\end{remark}

\subsection{The CCR is visible in the dynamics}
\label{subsec:CCR_in_dynamics}

\begin{proposition}[Pairwise dynamical noncommutativity]
\label{prop:CJ_commutator}
Extend $\log$ and $\Om$ additively to positive rationals and define $W_r$ by
\eqref{eq:Wn_definition} for $r\in\Qp$.  For $m,n\in\N$, put
\begin{equation}
\Phi_{m,n}
=
\vartheta\bigl(\Om(m)\log n-\log m\,\Om(n)\bigr).
\label{eq:Phi_mn}
\end{equation}
Then
\begin{align}
[C_m,J_n]
=
\frac12\sin\!\left(\frac{\Phi_{m,n}}{2}\right)
\Bigl(
W_{mn}+W_{m/n}-W_{n/m}-W_{1/(mn)}
\Bigr).
\label{eq:Cm_Jn_commutator}
\end{align}
\end{proposition}

\begin{proof}
Expand $C_m$ and $J_n$ using
\eqref{eq:Cn_definition}--\eqref{eq:Jn_definition}.  Apply
\eqref{eq:weyl_commutator} to the four pairs
$(z_m,z_n)$, $(z_m,-z_n)$, $(-z_m,z_n)$, and $(-z_m,-z_n)$.
\end{proof}

\begin{corollary}[Second-order arithmetic force]
\label{cor:second_order_force}
At $t=0$,
\begin{align}
\ddot Q(0)
=
\frac{iE_0^2}{2\hbar^2}
\sum_{m,n\ge2}
&\frac{\mu(m)\mu(n)x_n}{(mn)^{\sigma}}
\sin\!\left(\frac{\Phi_{m,n}}{2}\right)
\notag\\
&\times
\Bigl(
W_{mn}+W_{m/n}-W_{n/m}-W_{1/(mn)}
\Bigr),
\label{eq:Q_second_derivative}
\end{align}
and the analogous momentum equation is
\begin{align}
\ddot P(0)
=
\frac{iE_0^2}{2\hbar^2}
\sum_{m,n\ge2}
&\frac{\mu(m)\mu(n)y_n}{(mn)^{\sigma}}
\sin\!\left(\frac{\Phi_{m,n}}{2}\right)
\notag\\
&\times
\Bigl(
W_{mn}+W_{m/n}-W_{n/m}-W_{1/(mn)}
\Bigr).
\label{eq:P_second_derivative}
\end{align}
The double sums converge in operator norm for $\sigma>1$.
\end{corollary}

\begin{proof}
Differentiate \eqref{eq:Q_equation_motion} and
\eqref{eq:P_equation_motion} once more and use
Proposition~\ref{prop:CJ_commutator}.  Polynomial logarithmic weights are
summable against $(mn)^{-\sigma}$ for $\sigma>1$.
\end{proof}

\begin{corollary}[The interaction disappears with the Weyl cocycle]
\label{cor:interaction_disappears}
If $\vartheta=0$, then
\begin{equation}
\Phi_{m,n}=0
\label{eq:Phi_zero}
\end{equation}
for all $m,n$, and all pairwise terms in
\eqref{eq:Q_second_derivative}--\eqref{eq:P_second_derivative} vanish.
Thus the second-order cross-arithmetic force displayed above is generated
entirely by the exact Heisenberg cocycle.
\end{corollary}

\section{Exact non-localization inside the arithmetic dynamics}
\label{sec:nonlocalization}

\begin{proposition}[Schwartz invariance and distributional extension]
\label{prop:schwartz_distribution_extension}
For $\sigma>1$, the Hamiltonian $H_{\mu,\sigma,\vartheta}$ maps
$\Sch(\R)$ continuously into itself.  Consequently its coordinate formula
\eqref{eq:hamiltonian_coordinate_action} extends termwise to tempered
distributions.  For every fixed $T\in\Sch'(\R)$ the resulting series
converges in $\Sch'(\R)$.
\end{proposition}

\begin{proof}
For every Schwartz seminorm $p_{a,b}(\phi)=\sup_q|q^a\phi^{(b)}(q)|$,
the coordinate action \eqref{eq:weyl_coordinate_action} gives a bound of the
form
\begin{equation}
p_{a,b}(W_n\phi)
\le
C_{a,b}\bigl(1+|x_n|+|y_n|\bigr)^{N_{a,b}}
\sum_{j\le b}p_{a+j,b-j}(\phi),
\label{eq:weyl_schwartz_bound}
\end{equation}
for some finite integer $N_{a,b}$.  Since $x_n$ and $y_n$ grow at most
logarithmically and
\begin{equation}
\sum_{n\ge2}n^{-\sigma}(1+\log n)^N<\infty
\qquad(\sigma>1)
\label{eq:log_weight_summability}
\end{equation}
for every fixed $N$, the Hamiltonian series converges in every Schwartz
seminorm.  Translation and multiplication by a smooth phase act continuously
on $\Sch'$; the same logarithmic summability yields convergence after
applying the series to any fixed tempered distribution.
\end{proof}

\begin{proposition}[No exact position point is a stationary orbital]
\label{prop:no_delta_eigenstate}
Let $\sigma>1$, $\ell>0$, and $E_0>0$.  For no $q_0\in\R$ and no scalar
$E\in\C$ does the tempered distribution $\delta_{q_0}$ satisfy
\begin{equation}
H_{\mu,\sigma,\vartheta}\delta_{q_0}
=E\delta_{q_0}.
\label{eq:delta_eigenstate_impossible}
\end{equation}
\end{proposition}

\begin{proof}
By \eqref{eq:weyl_coordinate_action}, $W_n\delta_{q_0}$ is a nonzero scalar
multiple of $\delta_{q_0+x_n}$ and $W_n^*\delta_{q_0}$ is a nonzero scalar
multiple of $\delta_{q_0-x_n}$.  In particular, the $n=2$ term of
\eqref{eq:mobius_hamiltonian} contributes a nonzero multiple of
\begin{equation}
\delta_{q_0+\ell\log2}
\label{eq:delta_shift_plus_two}
\end{equation}
and a nonzero multiple of
\begin{equation}
\delta_{q_0-\ell\log2}.
\label{eq:delta_shift_minus_two}
\end{equation}
No other positive integer has the same positive displacement
$\ell\log2$, and no positive displacement equals its negative.  Hence these
support points cannot cancel against other terms or against $E\delta_{q_0}$.
Equation~\eqref{eq:delta_eigenstate_impossible} is impossible.
\end{proof}

\begin{remark}[Relation to exact Heisenberg impossibility]
\label{rem:nonlocalization_relation}
Theorem~\ref{thm:no_joint_eigenstate} is a universal consequence of the CCR.
Proposition~\ref{prop:no_delta_eigenstate} is more specific: the M\"obius
Hamiltonian itself does not admit an exactly position-localized stationary
mode.  The reason is not a measurement uncertainty but the exact translation
content of the arithmetic Weyl operators.
\end{remark}

\section{What is genuinely new in Version 2 and what remains standard}
\label{sec:status_classification}

\subsection{Standard ingredients}
\label{subsec:standard_ingredients}

The following ingredients are standard and are not presented as new:

\begin{enumerate}[leftmargin=*]
\item the Schr\"odinger representation of $Q$ and $P$ and the CCR
\eqref{eq:CCR};
\item Weyl operators and the Weyl product
\eqref{eq:weyl_product};
\item the arithmetic functions $\mu$, $\Om$, and classical M\"obius
inversion;
\item weighted $\ell^1$ Banach-algebra estimates and Neumann inversion.
\end{enumerate}

\subsection{Original notebook constructions}
\label{subsec:original_notebook_constructions}

The following are the constructions/results introduced in this notebook and
are the candidates for substantive originality:

\begin{enumerate}[leftmargin=*]
\item the arithmetic phase-space embedding
$n\mapsto(\ell\log n,\kappa\Om(n))$ coupled to the Weyl representation;
\item the cocycle \eqref{eq:chi_definition} as a deformation of Dirichlet
convolution;
\item the two-sided Heisenberg-twisted M\"obius inverse
$z_{\vartheta}$;
\item the exact prime-ray/prime-pair law
\eqref{eq:ztheta_prime_power}--\eqref{eq:ztheta_semiprime_defect};
\item the evenness theorem \eqref{eq:ztheta_evenness};
\item the M\"obius Hamiltonian \eqref{eq:mobius_hamiltonian} and its exact
position--momentum equations
\eqref{eq:Q_equation_motion}--\eqref{eq:P_equation_motion};
\item the pairwise second-order force controlled by
$\sin(\Phi_{m,n}/2)$.
\end{enumerate}

\begin{warning}[Priority]
\label{warn:priority}
``Introduced in this notebook'' is not a bibliographic priority claim.  The
present version deliberately develops the mathematics first.  A separate
literature audit must be performed before describing any item above as new in
the literature.
\end{warning}

\section{Open problems}
\label{sec:open_problems}

\begin{question}[Analytic continuation]
\label{q:analytic_continuation}
Does $\zeta_{\mathrm H}(s;\vartheta)$ admit a natural meromorphic continuation
from $\Re s\ge2$ to a larger domain for fixed nonzero $\vartheta$?
\end{question}

\begin{question}[Functional symmetry]
\label{q:functional_symmetry}
Is there a completion of $\zeta_{\mathrm H}(s;\vartheta)$ with a symmetry
that reduces at $\vartheta=0$ to the classical reflection $s\mapsto1-s$?
A nontrivial answer would be required before using the deformation to say
anything substantive about the critical line.
\end{question}

\begin{question}[Spectral zeros]
\label{q:spectral_zeros}
Can the zero set of a continued $\zeta_{\mathrm H}$ be realized as a spectral
or resonance set of an operator built intrinsically from
$H_{\mu,\sigma,\vartheta}$, rather than as zeros inserted into a multiplier?
\end{question}

\begin{question}[Classical-zero rigidity]
\label{q:classical_zero_rigidity}
If an analytic two-variable continuation through $\vartheta=0$ exists, does
the exact evenness in $\vartheta$ impose any second-order constraint on the
motion of the classical Riemann zeros beyond
Proposition~\ref{prop:conditional_zero_rigidity}?
\end{question}

\begin{question}[Choice of the second additive coordinate]
\label{q:additive_coordinate}
Which completely additive arithmetic functions $A(n)$, replacing $\Om(n)$
in $y_n=\kappa A(n)$, produce inequivalent deformations?  The choice
$A=\Om$ is distinguished here by its direct relation to the sign of $\mu$ on
square-free integers, but a classification is open.
\end{question}

\section{Notebook history}
\label{sec:notebook_history}

\begin{remark}
\label{rem:history_v1}
Version 1 used M\"obius-weighted logarithmic translations, Fourier
multipliers, an eta multiplier in the critical strip, generalized plane-wave
null modes, and a standard CCR non-localization statement.  Its main defect
was that the arithmetic multiplier and the Heisenberg pair were juxtaposed
rather than coupled.
\end{remark}

\begin{remark}
\label{rem:history_v2}
Version 2 removes the critical-strip multiplier as the central mechanism and
replaces ordinary Dirichlet convolution by the Weyl-cocycle product
\eqref{eq:twisted_convolution}.  Heisenberg noncommutativity now changes the
M\"obius inverse coefficients themselves and enters the exact equations of
motion of the M\"obius Hamiltonian.
\end{remark}

\begin{thebibliography}{99}

\bibitem{Heisenberg1927}
W.~Heisenberg,
\emph{\"Uber den anschaulichen Inhalt der quantentheoretischen Kinematik und
Mechanik},
Zeitschrift f\"ur Physik \textbf{43} (1927), 172--198.
doi:10.1007/BF01397280.

\bibitem{Folland1989}
G.~B.~Folland,
\emph{Harmonic Analysis in Phase Space},
Annals of Mathematics Studies 122,
Princeton University Press, Princeton, 1989.

\bibitem{ReedSimon1980}
M.~Reed and B.~Simon,
\emph{Methods of Modern Mathematical Physics, Vol.~I: Functional Analysis},
revised and enlarged ed., Academic Press, 1980.

\bibitem{Apostol1976}
T.~M.~Apostol,
\emph{Introduction to Analytic Number Theory},
Springer, New York, 1976.
doi:10.1007/978-1-4757-5579-4.

\end{thebibliography}

\end{document}
